% Generated draft: requires mathematical and source review.
\documentclass[11pt]{article}
\usepackage[margin=25mm]{geometry}
\usepackage{fontspec}
\setmainfont{Latin Modern Roman}
\usepackage{amsmath,amssymb,mathtools,unicode-math}
\setmathfont{Latin Modern Math}
\usepackage{enumitem,needspace,xurl,hyperref}
\setlength{\parindent}{0pt}
\setlength{\parskip}{5pt}
\setlength{\emergencystretch}{3em}
\newcommand{\N}{\mathbb{N}}
\newcommand{\Z}{\mathbb{Z}}
\newcommand{\Q}{\mathbb{Q}}
\newcommand{\R}{\mathbb{R}}
\newcommand{\C}{\mathbb{C}}
\newcommand{\F}{\mathbb{F}}
\newcommand{\A}{\mathbb{A}}
\newcommand{\PP}{\mathbb{P}}
\newcommand{\E}{\mathbb{E}}
\newcommand{\eps}{\varepsilon}
\newcommand{\dd}{\,\mathrm{d}}
\begin{document}
\section*{UnsolvedMath ID: 700580}
\textbf{Source problem code:} catalogue entry
\subsection{Definitions and mathematical statement}
A topological space \(K\) is scattered if every nonempty subset \(A\subseteq K\), with its subspace topology, has an isolated point. For a topological group \(G\), write
\[
\mathcal S_{\mathrm{nc}}(G)=\{H\leq G: H\text{ is closed in }G\text{ and is not compact}\}.
\]
The source asks whether every scattered compact space \(K\) admits a topological embedding
\[
e:K\longrightarrow \bigl(\mathcal S_{\mathrm{nc}}(G),\tau_E\bigr)
\]
for some locally compact group \(G\), which may depend on \(K\). Here a topological embedding is a continuous injection that is a homeomorphism onto its image with the subspace topology, and \(\tau_E\) denotes the source's \(E\)-topology.

The supplied record does not define the \(E\)-topology, so the topology on the target remains unresolved. It also does not explicitly state the separation conventions for ``compact'' or ``locally compact group,'' including whether Hausdorffness is required. These conventions must be recovered from the source before this becomes a fully specified mathematical statement.
\subsection{Short English statement}
Does every scattered compact space occur, up to homeomorphism, as a subspace of the space of closed noncompact subgroups of some locally compact group, equipped with the source's \(E\)-topology? The definition of that topology and the separation conventions remain unresolved.
\subsection{Sources}
\begin{itemize}
\item Local UnsolvedMath registry, supplied record with permanent ID \(700580\), exact source code \texttt{catalogue entry}, title ``Kourovka Notebook Problem 9.47.'' Attribution: I. V. Protasov, proposed in 1984.
\item \emph{The Kourovka Notebook (Unsolved Problems in Group Theory)}, 9th issue (1984), Problem 9.47; also identified in the supplied record with the 21st issue edition (2026), \texttt{arXiv:1401.0300v48}. The supplied collection description names E. I. Khukhro and V. D. Mazurov as editors.
\item The supplied background reports an explicitly unverified AI claim of a negative answer, attributed to unverified AI solutions: M. Kinyon and J. Urban, preprint of 17 September 2026, \url{https://kourovkanotebookorg.wordpress.com/wp-content/uploads/2026/09/kourovka-mathematics.pdf}.
\item The same explicitly unverified claim also cites Achyuth Jayadevan, preprint of 18 September 2026, \url{https://zenodo.org/records/22838025}.
\item Supplied status curation dated 2026-10-16 retains the canonical label ``open,'' classifies the evidence as \texttt{unrefereed\_claim}, and reports a \(\#\) marker for an unverified AI solution claim in the October 2026 21st issue. This is supplied dated curation, not an independently verified literature assessment. Its status-review quotation is truncated; the background supplies the complete links reproduced above.
\end{itemize}
\end{document}
